\documentclass[11pt]{article}
\usepackage{fontspec}
\setmainfont{Times New Roman}
\setsansfont{Arial}
\setmonofont{Consolas}
\usepackage{amsmath,amssymb}
\usepackage{geometry}
\usepackage{microtype}
\geometry{a4paper,margin=3cm}
\setlength{\parindent}{1.25em}
\setlength{\parskip}{0pt}
\emergencystretch=4em

\begin{document}

\begin{center}
{\Large\bfseries 2. Konstrukcija sistema form za ternarnu\\ bikvadratičnu formu}\par
\vspace{1.5em}
\emph{Journal für die reine und angewandte Mathematik} 134 (1908), s. 23--90 i dvě tabely
\end{center}

\vspace{3em}

\begin{center}
\textbf{Vvod.}
\end{center}

Sistem form ternarnoj bikvadratičnoj formy izslědovali Gordan, Maisano i Pascal.\footnote{Kratky izvod iz vvoda i glavy I pojavil se v \emph{Sitzungsberichte der physikalisch-medizinischen Sozietät Erlangen 1907}, s. 176--179.}\footnote{P. Gordan, \emph{Über das volle Formensystem der ternären biquadratischen Form}
\[
f=x_1^3x_2+x_2^3x_3+x_3^3x_1.
\]
(\emph{Math. Annalen}, tom XVII (1880), s. 217--233.)
G. Maisano, 1) \emph{Sistemi completi dei primi cinque gradi della forma ternaria biquadratica e degl' invarianti, covarianti e contravarianti di sesto grado}. (Giorn. di Battaglini XIX.)
2) \emph{Sui covarianti indipendenti di $6^\circ$ grado nei coefficienti della forma biquadratica ternaria}. (Rend. Circ. Mat. di Palermo I, 1887.)
E. Pascal, \emph{Contributo alla teoria della forma ternaria biquadratica e delle sue varie decomposizioni in fattori}. (Memoria premiata dalla R. Accademia delle scienze fisiche e matematiche di Napoli. 1905.)}
Gordan konstruuje polny sistem form, sostavjeny iz 54 konstrukcij, za specialnu avtomorfnu formu
\[
f=x_1^3x_2+x_2^3x_3+x_3^3x_1,
\]
na osnově podobnyh principov, ktore on dal za sistemy form v binarnoj oblasti.

U Maisana za obću bikvadratičnu formu konstruovane sut vse formy do petogo reda uključiteljno,\footnote{Pod "redom" tuto razuměje se stepenj v koeficientah, a pod "stepenem" stepenj v prěměnnyh.}
a takože něktore invarianty, kovarianty i kontravarianty vyšego reda. On koristaje metodu, ktoru Gordan priměnil v prvom tomu \emph{Mathematische Annalen} k ternarnoj kubičnoj formě. Pascal, opirajųći se na rezultaty Maisana, glavno se zanima pytanjem razklada bikvadratičnoj formy na faktory.\footnote{V svojoj drugoj kratkoj notě Maisano probuje dokazati linearnu zavisimost trěh kovariantov reda 6 i stepena 6. Nasuprot tomu ne vpolně dokončenomu dokazu Pascal mni, že dokazal linearnu nezavisimost tych trěh kovariantov reda 6, a takože trěh invariantov reda 9. No to, že linearno odnošenje v istině obstaje medžu trěmi kovariantami, s jednoj strany, i trěmi invariantami, s drugoj strany, pokazujut formuly (13), § 11, i (22), primětka k § 17, tutoj raboty, kde ta odnošenja sut dana javně. Po soobćenju Pascala protivorěčje objasnjaje se čislovoj pogrěškoju v izrazu jego kontravarianta $p$ (tuto označenogo $\varrho$), s. 46 jego statji.}

{\itshape Cilj slědujućih izslědovanj jest konstrukcija sistema form za obću ternarnu bikvadratičnu formu. V toj rabotě dajut se samo glavne osnovy i konstruuje se tak nazyvany "odnosno polny sistem".\footnote{Za termin gledati § 3a.}}

Rabota těsno primykaje k rabotě Gordana; vendar principy, ktore tam byli dane samo v obćem vidu, trěbalo bylo najprvo izrabotati v podrobnostjah. S pomočju teoremy o svęzi medžu svijanjami i črez vvedenje "ręda form" (§ 1) redukcijske teoremy formulujut se točno i doplnjajut se (§ 3), dokolě rekursivna konstrukcija specialnyh razvitij v rędy (§ 2) daje računno sredstvo za stvarno provedenje redukcij.

Osnovna myslj konstrukcije sistemov ternarnyh form jest taka sama kako v binarnoj oblasti. Počinajući od prvogo odnosno polnogo sistema --- sistema form prěnesenyh iz binarnogo slučaja, --- po opreděljenom zakonu prěhodi se k sistemam s vse vyšim modulom. Ta postupnost ide dotud, dokud sistem někojego modula, ako sam modul vzęti za osnovnu formu, ne stane konečnym i poznanym, ili dokud modul ne možno svesti k formam, ktore imajut invarianty kako faktory. V prvom slučaju absolutno polny sistem nastaje transvekcijeju odnosno polnogo sistema so sistemom modula; v drugom slučaju odnosno polny i absolutno polny sistem sut identične. Iz obćego Hilbertovogo dokaza konečnosti sistemov form slěduje, že ta procedura neobhodno mora dojti do konca.

V našem slučaju modul $(abc)$ prvogo odnosno polnogo sistema (§ 4) svodi se k modulam
\[
\Delta=(abc)^2a_x^2b_x^2c_x^2\quad\text{i}\quad \nu=(abu)^4
\]
(§ 5), a potom nahodi se odnosno polny sistem modulo $\nu$ (§ 6). Ty rezultaty sut uže dane v rabotě Gordana za obću bikvadratičnu formu; naš sposob izvoda jednak lehče prěnosi se na formy vyšego stepena.

Kako ręd modulov teper vybirajemo formy (§ 7):
\[
\nu,\qquad \nu(\nu)=(\nu\nu_1x)^4=s_x^4,\qquad \nu(s)=(ss'u)^4,\ldots .
\]

Pri konstrukciji odnosno polnogo sistema modulo $s$ dvě kvadratične formy, nastavajuće v sistemu, $u_\varrho^2$ i $t_x^2$, pridajut se kako moduly (§§ 9 i 10), i odtud konstruuje se odnosno polny sistem modulo $(s,\varrho,t)$. Potom pokazuje se (§ 17), že modul $(ss'u)^4$ jest reducibilny k modulam $\varrho$ i $t$. Črez to slědujući vyši sistem prěhodi v odnosno polny sistem modulo $(\varrho,t)$. No simultanny sistem dvuh kvadratičnyh form jest konečny i poznany, a v obćem slučaju sostavjen jest iz 20 konstrukcij;\footnote{Srav. Clebsch--Lindemann, \emph{Vorlesungen über Geometrie} (tom I, s. 288 i dalje), sistem dvuh kogredientnyh form. Gordan, \emph{Über Büschel von Kegelschnitten} (\emph{Math. Ann.}, tom XIX, s. 530), sistem dvuh kontragredientnyh form.}
tym s konstrukcijeju odnosno polnogo sistema modulo $(\varrho,t)$ dosegnuty jest vyše označeny konec. Transvekcija togo sistema so sistemom $(\varrho,t)$ za konstrukciju absolutno polnogo sistema ostaje zadržana za pozdnějši čas.

Kako odnosno polny sistem modulo $(\varrho,t)$ najdene sut 331 konstrukcije; v prilagajenej tabeli one sut uporędkovane po stepenu v prěměnnyh $x$ i $u$.

Na koncu dajemo pregled vvedenyh simbolov:
\begin{align*}
f&=a_x^4,\\
\theta&=\theta_x^4u_\vartheta^2=(abu)^2a_x^2b_x^2,\\
K&=K_x^6u_k^3=(a\theta u)u_\vartheta^2a_x^3\theta_x^3,\\
j&=u_j^6=(a\theta u)^4u_\vartheta^2,\\
\Delta&=\Delta_x^6=a_\vartheta^2a_x^2\theta_x^4=(abc)^2a_x^2b_x^2c_x^2,\\
N&=N_x^8u_\eta=(a\Delta u)a_x^3\Delta_x^5,\\
\nu&=u_\nu^4=(abu)^4,\\
H&=H_x^2u_\eta^4=(\nu\nu_1x)^2u_\nu^2u_{\nu_1}^2,\\
L&=L_x^3u_l^6=(\nu\eta x)H_x^2u_\nu^3u_\eta^3,\\
g&=g_x^6=(\nu\eta x)^4H_x^2,\\
\sigma&=u_\sigma^6=H_\nu^2u_\nu^2u_\eta^4,\\
s&=s_x^4=(\nu\nu_1x)^4,\\
Z&=Z_x^4u_\zeta^2=(ss'u)^2s_x^2s_x'^2,\\
\varrho&=u_\varrho^2=\theta_\nu^4u_\vartheta^2,\\
t&=t_x^2=a_\eta^4H_x^2,\\
i&=a_\nu^4,\\
J&=s_\nu^4.
\end{align*}

\end{document}
